% Section 3. Sources: 607_bounds_paper/paper.tex section 3 (weight inequality, four points, best weight, ledger);
% 641_tape_window_bounds/PROOF.md Lemma U; 607_bounds_paper/paper_word.tex section 3.1 (edge certificate);
% 1003_paired_path_consistency/PROOF.md section 1 (finite SOS theorem), section 4 (field is not a potential);
% 1251_referee_upper_12297/STATEMENT.md and FINAL.md section 2 (normalization); 1201_upper_structures/FINAL.md
% (edge-class inventory: Perron vectors give equality).
%%%%%%%%%%%%%%%%%%%%%%%%%%%%%%%%%%%%%%%%%%%%%%%%%%%%%%%%%%%%%%%%%%%%%%%%
\section{From weights to an upper bound}\label{sec:weights}
%%%%%%%%%%%%%%%%%%%%%%%%%%%%%%%%%%%%%%%%%%%%%%%%%%%%%%%%%%%%%%%%%%%%%%%%

A \emph{history} is a walk of the tape machine that starts at the root and may end anywhere. A \emph{successor} of a
tape is a tape one move later; a \emph{predecessor} is a tape one move earlier. By Theorem~\ref{thm:tape} and
\eqref{eq:fekete}, bounding $\mu$ from above means bounding the number of closed walks.

\subsection{One weight}

Let $W$ assign a positive number $W(T)$ to every tape $T$ of $\cls$, and suppose that for some $R$ and every tape $T$
of $\cls$,
\begin{equation}\label{eq:sub}
  \sum_{T'\ \text{successor of}\ T} W(T') \;\le\; R\,W(T).
\end{equation}
Intuitively, every tape has at most $R$ continuations, when continuations are counted by weight and not by number.

\begin{lemma}\label{lem:oneweight}
If \eqref{eq:sub} holds at every tape of $\cls$, then $M_n\le R^{2n-2}$ for every $n$, and $\mu\le R^2$.
\end{lemma}

\begin{proof}
Give each history the weight of its last tape. One more move replaces a history ending at $T$ by one history for each
successor of $T$, and by \eqref{eq:sub} their weights add up to at most $R\,W(T)$. So the histories of $t$ moves have
total weight at most $R^tW(\text{root})$. Those that end at the root weigh $W(\text{root})$ each, and all others have
positive weight, so at most $R^t$ histories of $t$ moves end at the root. Put $t=2n-2$ and use
Theorem~\ref{thm:tape}.
\end{proof}

Intuitively, the proof keeps a running total of the weights of all histories, and this total grows by at most a
factor $R$ per move. After four moves eight histories are back at the root, the eight meanders of order $3$; the
others are counted too, which makes the bound loose but not wrong. Four points of the argument are easy to get
wrong~\cite{BoundsPaper}. The inequality is about the \emph{sum}: a single successor may weigh more than the present
tape. A history carries \emph{one} weight, that of its current tape; the factors $W(T')/W(T)$ along a history
telescope. Histories that do not return to the root are \emph{kept}, with positive weight. And weights are \emph{not
probabilities}: their only job is to satisfy \eqref{eq:sub}.

With all weights equal, \eqref{eq:sub} holds with $R=4$, the largest number of moves at a tape, and gives the trivial
bound $\mu\le16$. Unequal weights help because tapes differ in how many continuations they have; for example a tape
with the head on a $-$ has no merge, and a tape with the head at an end has one head move fewer. The ideal weight would
be the number of ways to return to the root in a given number of moves, which is the original counting problem. The
weights used below are tractable stand-ins, and they are valid only because the inequality is checked everywhere. In
matrix language, \eqref{eq:sub} says $AW\le RW$ for the $0/1$ matrix $A$ of the moves, and
Lemma~\ref{lem:oneweight} is the easy half of the Collatz--Wielandt characterization of the Perron root, the largest
eigenvalue of a nonnegative matrix~\cite{Seneta}.


\subsection{Two weights and reciprocal edge factors}

The certificate of Section~\ref{sec:upper} uses a sharper test. It looks at each tape both forward, along its moves,
and backward, along the moves that lead into it, and it lets each move carry a factor that its reverse pays back.
Intuitively, a move may borrow weight as long as the reverse move returns it.

Let $X$ be a finite set of tapes and $A$ its $0/1$ matrix of moves, with rows and columns indexed by the tapes:
$A_{TT'}=1$ if a move goes from $T$ to $T'$, both in $X$. Choose positive weights $W(T)$ and $Z(T)$ on $X$. Call
$\{T,T'\}$ an \emph{edge} if a move goes from $T$ to $T'$ or from $T'$ to $T$, and choose for every edge a positive
factor $\theta_{TT'}$ with
\begin{equation}\label{eq:recip}
  \theta_{T'T}=1/\theta_{TT'}.
\end{equation}
The \emph{row value} of tape $T$ is
\begin{equation}\label{eq:row}
  \phi_T \;=\; \frac12\sum_{T':\ T\to T'}\theta_{TT'}\,\frac{W(T')}{W(T)}\;+\;\frac12\sum_{T':\ T'\to T}\theta_{TT'}\,\frac{Z(T')}{Z(T)}.
\end{equation}
The first sum runs over successors, the second over predecessors. In both, the factor is the one oriented from the
present tape $T$ toward its neighbor $T'$, not the one in the direction of the move. Confusing the two orientations
would invalidate the certificate.

\begin{lemma}[paired rows]\label{lem:paired}
If $\phi_T\le R$ for every $T\in X$, then the spectral radius of $A$ is at most $R$.
\end{lemma}

\begin{proof}
Put
\[ B_{TT'}=A_{TT'}\sqrt{\frac{Z(T)\,W(T')}{W(T)\,Z(T')}},\qquad
   \vartheta_{TT'}=\theta_{TT'}\sqrt{\frac{W(T')\,Z(T')}{W(T)\,Z(T)}} . \]
The matrix $B$ is obtained from $A$ by conjugating with a positive diagonal matrix, so it has the same spectral
radius $\rho(A)$. Let $\widehat B=\frac12(B+B^{\mathsf T})$ be its symmetric part. Then
$\vartheta_{T'T}=1/\vartheta_{TT'}$, and a short computation gives $\widehat B_{TT'}\vartheta_{TT'}=\frac12\bigl(A_{TT'}W(T')/W(T)+A_{T'T}Z(T')/Z(T)\bigr)\theta_{TT'}$,
so that $\phi_T=\sum_{T'}\widehat B_{TT'}\vartheta_{TT'}$. For every edge and every real vector $v$,
$2v_Tv_{T'}\le \vartheta_{TT'}v_T^2+\vartheta_{T'T}v_{T'}^2$, since the difference is
$\bigl(\sqrt{\vartheta_{TT'}}\,v_T-v_{T'}/\sqrt{\vartheta_{TT'}}\,\bigr)^2$. Summing over the edges with the weights
$\widehat B_{TT'}\ge0$ gives
\[ v^{\mathsf T}\widehat Bv\;\le\;\sum_T\phi_Tv_T^2\;\le\;R\,\lvert v\rvert^2 . \]
A move always changes the tape, so $A$ has no diagonal entries. The matrix $B\ge0$ has a nonnegative eigenvector $v$
for its spectral radius $\rho(A)$, and for it $v^{\mathsf T}\widehat Bv=v^{\mathsf T}Bv=\rho(A)\lvert v\rvert^2$. Hence
$\rho(A)\le R$.
\end{proof}

\begin{corollary}\label{cor:paired}
If $\phi_T\le R$ holds for every finite set $X$ of tapes of $\cls$ and every $T\in X$, then $M_n\le R^{2n-2}$ for
every $n$, and $\mu\le R^2$.
\end{corollary}

\begin{proof}
Fix $n$ and $t=2n-2$, and let $X$ be the tapes of $\cls$ with at most $t+3$ letters; every closed walk of $t$ moves
stays in $X$, and so does every concatenation of such walks. So $(A^{\ell t})_{\text{root},\text{root}}\ge
M_n^{\,\ell}$ for every $\ell\ge1$. Taking $\ell t$-th roots and letting $\ell\to\infty$ gives
$M_n^{1/t}\le\rho(A)\le R$.
\end{proof}

The argument uses finite matrices only; it never assumes that a change of weights is valid on the infinite set of
all tapes.
Restricting to $X$ only deletes positive terms from \eqref{eq:row}, so in practice one checks the row values with all
neighbors present.

\paragraph{Why two weights and edge factors help.}
Lemma~\ref{lem:oneweight} looks at successors only. Averaging the forward and the backward view lets the slack of
one direction (the amount by which its sum falls below $R$) pay for the other. The edge factors add a freedom that no
weight on tapes has.
The ratios $W(T')/W(T)$ of a single weight multiply to one around every closed cycle of moves; reciprocal edge factors need
only invert along the \emph{same} edge, so around a cycle they can multiply to anything. They are parameters of the
inequality, not move probabilities. On a finite set of tapes, taking for $W$ and $Z$ the right and left Perron
vectors of $A$ (its eigenvectors for the spectral radius) and all factors equal to one makes every row value equal to
the spectral radius. So the Perron vectors are the ideal choice, and improving a certificate means describing them
better by features that can be read near the head.

It is natural to take $Z$ to be $W$ of the complemented tape $\cpl T$, in which every letter is changed. Reversing a
move and complementing the tape gives a forward move: the reverse of the split $[+]\mapsto+[-]+$ is
$+[-]+\mapsto[+]$, whose complement $-[+]-\mapsto[-]$ is a merge, and the same holds for the other splits and merges;
head moves are their own reverses. So the predecessors of $T$ are the complements of the successors of $\cpl T$, and
the backward sum of \eqref{eq:row} is a forward sum for the complemented tape. The missing merge at $+[+]+$ has no
edge in either direction.
